\section{评测体系与质检闭环}
\label{sec:eval}

流水线的闸门保证交付物“齐全、新鲜、数字有出处”，但回答不了“这份研报好不好”。本节介绍我们的评分标准、计分方式、盲评设置、多轮质检闭环，以及如何校准评审模型本身。

\subsection{评分标准的来源}
评分标准由两条思路合成。\textbf{思路一}从“卖方深度报告在回答哪些问题”出发：按常规体例逐部分列出读者要的答案——首页要能复述“为什么值这个价”，核心观点要说明与市场的分歧在哪里，驱动拆解要有“外部驱动 → 量价结构 → 收入利润”的定量链条，风险要落到具体假设——每个问题改写成一条可二元判定的条目，并写清判定标准与反例。权重按对结论的影响分三档：决定估值的环节 5 分，支撑判断的关键论证 4 分，完整性要求 2～3 分。\textbf{思路二}从工作流出发：把“年报 → 抽取 → 真值表 → 经营数据 → 预测 → 三表 → DCF → 文稿 → 排版”的每一次交接看作一个闭合点，上游的数在下游必须找得到、对得上、口径一致；这些条目几乎都能由程序判定，直接复用已有的校验工具。

\begin{table}[t]
\centering
\caption{评分标准的层级结构（v2）。L0 不计分但失败即封顶；N 为负向条目，命中扣分。}
\label{tab:rubric}
\small
\begin{tabularx}{\linewidth}{@{}p{1.9cm}cc p{1.6cm}X@{}}
\toprule
层 & 条目 & 分值 & 判定者 & 内容 \\
\midrule
L0 硬闸门 & 5 & — & 程序 & 数字护栏、三表勾稽、DCF 与三表同口径、免责声明、评级与估值一致；任一失败总分封顶 40 \\
L1 数据闭合 & 13 & 35 & 程序 & 抽取校验、分部预测基期闭合、股权价值桥加总、每股价值复算、情景排序、敏感性方向、终值占比、预测与历史衔接、正文估值与产物一致、外部信息可追溯、股价日期与来源、数字锁、人工复核清单 \\
L2 分析质量 & 12 & 45 & 评审模型 & 核心结论、预期差、商业模式、驱动拆解、历史业绩归因、行业与竞争、预测假设依据、估值方法与参数、估值交叉验证、情景与敏感性、风险落到假设、跟踪指标与失效条件 \\
L3 表达与合规 & 7 & 20 & 评审模型 & 故事线连贯、事实 / 推论 / 观点可区分、结论先行、无机器腔、数字密度、图表服务论点、假设与局限透明 \\
N 负向 & 4 & $-30$ & 程序 / 评审 & 无来源外部数据（$-8$，并封顶 70）、数字前后矛盾（$-8$，程序判）、聚合数构成描述错误（$-6$）、结论前后冲突（$-8$） \\
\bottomrule
\end{tabularx}
\end{table}

\subsection{计分}
所有计分都在 Python 中完成，评审模型只输出逐条的 $\{\text{id},\ \text{是否通过},\ \text{证据},\ \text{理由}\}$。设 $E_\ell$ 为第 $\ell$ 层得分，$D=\max\bigl(\sum_{i\in N,\,\text{命中}} p_i,\ D_{\min}\bigr)$ 为负向扣分（下限 $D_{\min}=-20$），$\mathcal{C}$ 为触发的封顶集合（L0 任一失败封顶 40；N-01 命中封顶 70），则
\begin{equation}
  S = \min\Bigl(\max\bigl(0,\ E_{L1}+E_{L2}+E_{L3}+D\bigr),\ \min\mathcal{C}\Bigr).
\end{equation}
缺判定的条目按未通过计（负向条目按未命中计）并单独列出。\textbf{交付标准}为：$S\ge 75$；L1、L2、L3 得分率分别不低于 80\%、60\%、60\%；且未命中任何负向条目。权重、封顶与门槛都写在评分标准文件中，修改后可直接重算，无需重跑评审。

评分标准经历了三个版本。v1 依据三家 v0 研报的人工质检调整：外部信息可追溯（L1-10）从“来源文件是否存在”改为核对引用编号与外部数字的对应；人工复核清单（L1-13）只认填写完整的事项；“无机器腔”（L3-04）只看正文中成为习惯写法的程序口吻，图表来源与免责声明不算；“数字密度”（L3-05）把数字连同其同比、占比算作一个数据点。v2 把“数字前后矛盾”改为程序判定，把数值以外的结论冲突拆为单独的 N-04。

\subsection{盲评}
评审在仓库与工作区之外单独建立的目录中进行，目录里只有交付物：研报 PDF、文稿、估值产物（DCF、预测摘要、分部预测、情景、非经营性资产、可比公司）、外部来源登记、人工复核清单，以及评审条目（只含编号、标题、判定标准与反例，\textbf{不含分值与层级权重}）。评审会话看不到写作与复核过程、模型名称和其他公司的结果；只允许读本目录、写 \file{judgments.json}，不联网。评审说明要求：每一条都给出判定；完全满足判定标准才算通过，\textbf{拿不准时判不通过}并写明缺什么；证据摘录原文并标注章节或页码。数据是否与年报一致、勾稽与估值算术是否正确已由程序核对，评审不必复算，只判断分析是否成立、论证是否充分、表达是否合格。

\subsection{质检闭环}
\label{sec:refine}

\begin{figure}[t]
  \centering
  \figslot{fig6_refine}{5.6cm}
  \caption{质检闭环与信息隔离。生成侧（工作区内）看不到评分标准与分值；评测侧（仓库与工作区之外）只看到交付物。不达标时，问题单按路由回传：数据与假设类回到分析师 / 主笔，措辞与结构类交给复核编辑，版面类交给排版设计师。每轮存档，结束时回滚到得分最高的一轮。}
  \label{fig:refine}
\end{figure}

\code{refine} 把生成、复核与质检串成多轮循环（算法 \ref{alg:refine}，图 \ref{fig:refine}）。每条评分条目都带一个路由字段（\code{generator} 或 \code{reviewer}）和对应阶段；未通过的条目被写成问题单，只含问题描述、理由与证据，不含条目原文和分值。路由回生成侧的问题会删除初稿备份与修改记录，使复核从头再来，避免旧的复核结论覆盖新的文稿。

\begin{algorithm}[t]
\caption{质检闭环（\code{refine}）}
\label{alg:refine}
\small
\begin{algorithmic}[1]
\Require 工作区 $W$，最大轮数 $K=3$，最小提升 $\Delta=2$
\For{$n=1,\dots,K$}
  \If{$n=1$} 依次运行五个会话 \Comment{已通过的阶段跳过}
  \Else
    \State 读取上一轮问题单 $I$，按路由分为 $I_{\text{gen}}, I_{\text{rev}}, I_{\text{layout}}$
    \If{$I_{\text{gen}}\ne\varnothing$} 删除初稿备份；修订会话处理 $I_{\text{gen}}$（需要时重跑估值与护栏）；随后复核与排版按闸门续做
    \Else{} 复核编辑处理 $I_{\text{rev}}$（数字锁照常）；排版设计师处理 $I_{\text{layout}}$
    \EndIf
  \EndIf
  \If{\code{status}($W$) 未全部通过} 停止，交给人 \EndIf
  \State 盲评并计分得 $S_n$，存档本轮文稿、复核文件与 PDF
  \If{达标 \textbf{or} $S_n-S_{n-1}<\Delta \wedge S_{n-1}-S_{n-2}<\Delta$ \textbf{or} 无可回传问题} \textbf{break} \EndIf
\EndFor
\State 回滚到 $\arg\max_n S_n$（若期间重跑过估值则只标注最佳轮次）
\State 若最佳轮仍未达标，把剩余问题追加进人工复核清单
\end{algorithmic}
\end{algorithm}

\paragraph{防止对评分标准过拟合。}多轮迭代天然有朝评分标准优化的风险（Goodhart 定律）。我们的对策有三：评分标准与完整评测结果放在工作区之外，生成方读不到；问题单只给理由与证据，不给条目原文、分值与通过线；L1 全部由程序判定，生成方无法通过措辞“说服”程序。设计中还预留了把评分标准拆成“反馈集”和“留出集”的方案：反馈集转成问题单，留出集只用于最终计分，两者得分差距过大即说明在应试。

\subsection{校准评审模型}
\label{sec:calib-method}
评审模型的判定只有在被检验之后才能用来驱动迭代。我们用三组证据回答三个不同的问题：

\begin{enumerate}
  \item \textbf{稳定性：同一份研报重复盲评，判定是否一致？}对同一交付物重复盲评 $k$ 次，计算逐条全一致的比例、两两不一致率与总分极差。
  \item \textbf{灵敏度：已知缺陷能否判出？}对一份合格研报的文稿施加确定性的“劣化算子”，制造已知缺陷后重新排版、重新盲评（表 \ref{tab:degrade}）。缺陷是程序造的，“该判不通过”是已知真值。同时统计非目标条目被连带改判的数量，并与重复评审本身的波动对比。原稿已被判出该问题的变体不计入命中率。
  \item \textbf{正确性：判出的问题，人是否认同？}人工质检页逐条复核评审判“不通过”的条目，计算判出问题的精确率。由于人工只复核了判不通过的条目，这组证据算不出漏判率，漏判由第 2 组补充。
\end{enumerate}

\begin{table}[t]
\centering
\caption{劣化算子。每个算子只修改文稿，确定性地制造一种已知缺陷。}
\label{tab:degrade}
\small
\begin{tabularx}{\linewidth}{@{}p{2.6cm}Xc@{}}
\toprule
变体 & 制造的缺陷 & 目标条目 \\
\midrule
contradiction & 只把摘要里的评级改成相反方向，正文不变 & N-02 \\
generic\_risks & 风险提示换成不对应任何假设的通用模板 & L2-11 \\
machine\_tone & 每两段注入一次程序口吻（如“根据程序计算结果”） & L3-04 \\
no\_cross\_check & 删掉相对估值与交叉验证小节 & L2-09 \\
no\_expectation\_gap & 删掉市场关注问题、预期差与“股价隐含”专题 & L2-02 \\
no\_tracking & 删掉跟踪指标与结论失效条件 & L2-12 \\
unsourced\_number & 行业部分插入一段无来源的市场规模与份额数字 & N-01 \\
\bottomrule
\end{tabularx}
\end{table}
